% GENERATED FILE. Edit canonical lessons, then run npm run generate:books.
% canonical-source-hash: fnv1a64:8758e20a67d26e66
% canonical-lessons: JA-C223-abunai, JA-C223-kiken, JA-C223-ikiru, JA-C223-tasukeru, JA-C223-kurushimu, JA-R223-review-222-223, JA-R223-review-final-223

\chapter{Danger and Staying Alive}
\label{ch:ja-danger-and-staying-alive}
\hlchaptermodality{223}

\begin{chapteropening}
\textbf{By the end of this chapter:} \emph{I can warn about danger and ask if someone is safe.}

Complete the last lesson of chapter 223: I can warn about danger and ask if someone is safe.
\end{chapteropening}

\section[\texorpdfstring{\ja{あぶない} (abunai)}{あぶない (abunai)}]{\ja{あぶない} (abunai) — dangerous}

\label{lesson:JA-C223-abunai}



\begin{quote}
Before the new one: say the Japanese for all right, then the Japanese for endurance.
\end{quote}

\subsection*{You'll want to know: \ja{あぶない}}
\textbf{\ja{あぶない}} — \emph{abunai} — \textquotedblleft{}dangerous\textquotedblright{}.

\begin{quote}
\textbf{\ja{この} \ja{みちは} \ja{あぶないですか。}} — \emph{kono michi wa abunai desu ka} — \textquotedblleft{}Is this road dangerous?\textquotedblright{}
\end{quote}

\begin{cousinweb}
\textbf{\ja{あぶない}} (\emph{abunai}) is a native Japanese word (wago), inherited rather than borrowed from Chinese; its kanji was attached later as a kun'yomi, and its older origin is not settled (JMdict, KANJIDIC2).
\end{cousinweb}

\begin{grammarlens}[title={What you've built}]
One more word for everyday life at A2.
\end{grammarlens}

\subsection*{Your turn}
\begin{itemize}
\raggedright
  \item \emph{Say it:} \emph{abunai}
  \item \emph{Say it:} \emph{abunai}, once more
  \item \emph{Say it:} \emph{kono michi wa abunai desu ka} — Is this road dangerous?
  \item \emph{Recall:} say the Japanese for all right, then the Japanese for endurance, then say \emph{abunai} again
\end{itemize}

\begin{tcolorbox}[breakable,colback=teal!4,colframe=teal!35!black,title={Before you move on}]
What does \ja{あぶない} mean? (\textquotedblleft{}Dangerous\textquotedblright{}.) Say it once more.
\end{tcolorbox}
\section[\texorpdfstring{\ja{きけん} (kiken)}{きけん (kiken)}]{\ja{きけん} (kiken) — danger}

\label{lesson:JA-C223-kiken}



\begin{quote}
Before the new one: say the Japanese for endurance, then the Japanese for dangerous.
\end{quote}

\subsection*{You'll want to know: \ja{きけん}}
\textbf{\ja{きけん}} — \emph{kiken} — \textquotedblleft{}danger\textquotedblright{}.

\begin{quote}
\textbf{\ja{いのちに} \ja{きけんは} \ja{ありません。}} — \emph{inochi ni kiken wa arimasen} — \textquotedblleft{}There is no danger to life.\textquotedblright{}
\end{quote}

\begin{cousinweb}
\textbf{\ja{きけん}} (\emph{kiken}) is Sino-Japanese (kango): two Chinese characters read with their on'yomi, readings borrowed from Chinese. They are \emph{ki}, \textquotedblleft{}danger\textquotedblright{} (Middle Chinese \emph{ngjwe}) and \emph{ken}, \textquotedblleft{}steep, risky\textquotedblright{} (\emph{xjemX}) (KANJIDIC2; Middle Chinese in Baxter's transcription).
\end{cousinweb}

\begin{grammarlens}[title={What you've built}]
One more word for everyday life at A2.
\end{grammarlens}

\subsection*{Your turn}
\begin{itemize}
\raggedright
  \item \emph{Say it:} \emph{kiken}
  \item \emph{Say it:} \emph{kiken}, once more
  \item \emph{Say it:} \emph{inochi ni kiken wa arimasen} — There is no danger to life.
  \item \emph{Recall:} say the Japanese for endurance, then the Japanese for dangerous, then say \emph{kiken} again
\end{itemize}

\begin{tcolorbox}[breakable,colback=teal!4,colframe=teal!35!black,title={Before you move on}]
What does \ja{きけん} mean? (\textquotedblleft{}Danger\textquotedblright{}.) Say it once more.
\end{tcolorbox}
\section[\texorpdfstring{\ja{いきる} (ikiru)}{いきる (ikiru)}]{\ja{いきる} (ikiru) — to be alive; to live (on)}

\label{lesson:JA-C223-ikiru}



\begin{quote}
Before the new one: say the Japanese for dangerous, then the Japanese for danger.
\end{quote}

\subsection*{You'll want to know: \ja{いきる}}
\textbf{\ja{いきる}} — \emph{ikiru} — \textquotedblleft{}to be alive; to live (on)\textquotedblright{}.

\begin{quote}
\textbf{\ja{さかなは} \ja{みずの} \ja{なかで} \ja{いきます。}} — \emph{sakana wa mizu no naka de ikimasu} — \textquotedblleft{}Fish live in water.\textquotedblright{}
\end{quote}

Polite \textbf{\ja{いきます}}, \emph{ikimasu}; an ichidan verb. It sounds the same as the polite form of \textbf{\ja{いく}}, \emph{iku}.

\begin{cousinweb}
\textbf{\ja{いきる}} (\emph{ikiru}) is a native Japanese word (wago), inherited rather than borrowed from Chinese; its kanji was attached later as a kun'yomi, and its older origin is not settled (JMdict, KANJIDIC2).
\end{cousinweb}

\begin{grammarlens}[title={What you've built}]
One more word for everyday life at A2.
\end{grammarlens}

\subsection*{Your turn}
\begin{itemize}
\raggedright
  \item \emph{Say it:} \emph{ikiru}
  \item \emph{Say it:} \emph{ikiru}, once more
  \item \emph{Say it:} \emph{sakana wa mizu no naka de ikimasu} — Fish live in water.
  \item \emph{Recall:} say the Japanese for dangerous, then the Japanese for danger, then say \emph{ikiru} again
\end{itemize}

\begin{tcolorbox}[breakable,colback=teal!4,colframe=teal!35!black,title={Before you move on}]
What does \ja{いきる} mean? (\textquotedblleft{}To be alive; to live (on)\textquotedblright{}.) Say it once more.
\end{tcolorbox}
\section[\texorpdfstring{\ja{たすける} (tasukeru)}{たすける (tasukeru)}]{\ja{たすける} (tasukeru) — to rescue; to save}

\label{lesson:JA-C223-tasukeru}



\begin{quote}
Before the new one: say the Japanese for danger, then the Japanese for to be alive.
\end{quote}

\subsection*{You'll want to know: \ja{たすける}}
\textbf{\ja{たすける}} — \emph{tasukeru} — \textquotedblleft{}to rescue; to save\textquotedblright{}.

\begin{quote}
\textbf{\ja{いしゃは} \ja{いのちを} \ja{たすけます。}} — \emph{isha wa inochi o tasukemasu} — \textquotedblleft{}Doctors save lives.\textquotedblright{}
\end{quote}

Polite \textbf{\ja{たすけます}}, \emph{tasukemasu}; an ichidan verb.

\begin{cousinweb}
\textbf{\ja{たすける}} (\emph{tasukeru}) is a native Japanese word (wago), inherited rather than borrowed from Chinese; its kanji was attached later as a kun'yomi, and its older origin is not settled (JMdict, KANJIDIC2).
\end{cousinweb}

\begin{grammarlens}[title={What you've built}]
One more word for everyday life at A2.
\end{grammarlens}

\subsection*{Your turn}
\begin{itemize}
\raggedright
  \item \emph{Say it:} \emph{tasukeru}
  \item \emph{Say it:} \emph{tasukeru}, once more
  \item \emph{Say it:} \emph{isha wa inochi o tasukemasu} — Doctors save lives.
  \item \emph{Recall:} say the Japanese for danger, then the Japanese for to be alive, then say \emph{tasukeru} again
\end{itemize}

\begin{tcolorbox}[breakable,colback=teal!4,colframe=teal!35!black,title={Before you move on}]
What does \ja{たすける} mean? (\textquotedblleft{}To rescue; to save\textquotedblright{}.) Say it once more.
\end{tcolorbox}
\section[\texorpdfstring{\ja{くるしむ} (kurushimu)}{くるしむ (kurushimu)}]{\ja{くるしむ} (kurushimu) — to suffer}

\label{lesson:JA-C223-kurushimu}



\begin{quote}
Before the new one: say the Japanese for to be alive, then the Japanese for to rescue.
\end{quote}

\subsection*{You'll want to know: \ja{くるしむ}}
\textbf{\ja{くるしむ}} — \emph{kurushimu} — \textquotedblleft{}to suffer\textquotedblright{}.

\begin{quote}
\textbf{\ja{かんじゃは} \ja{もう} \ja{くるしみません。}} — \emph{kanja wa mō kurushimimasen} — \textquotedblleft{}The patient is no longer suffering.\textquotedblright{}
\end{quote}

Polite \textbf{\ja{くるしみます}}, \emph{kurushimimasu}; a godan verb.

\begin{cousinweb}
\textbf{\ja{くるしむ}} (\emph{kurushimu}) is a native Japanese word (wago), inherited rather than borrowed from Chinese; its kanji was attached later as a kun'yomi, and its older origin is not settled (JMdict, KANJIDIC2).
\end{cousinweb}

\begin{grammarlens}[title={What you've built}]
One more word for everyday life at A2.
\end{grammarlens}

\subsection*{Your turn}
\begin{itemize}
\raggedright
  \item \emph{Say it:} \emph{kurushimu}
  \item \emph{Say it:} \emph{kurushimu}, once more
  \item \emph{Say it:} \emph{kanja wa mō kurushimimasen} — The patient is no longer suffering.
  \item \emph{Recall:} say the Japanese for to be alive, then the Japanese for to rescue, then say \emph{kurushimu} again
\end{itemize}

\begin{tcolorbox}[breakable,colback=teal!4,colframe=teal!35!black,title={Before you move on}]
What does \ja{くるしむ} mean? (\textquotedblleft{}To suffer\textquotedblright{}.) Say it once more.
\end{tcolorbox}
\section[(dialogue)]{Medicine and Safety - a review}

\label{lesson:JA-R223-review-222-223}



\begin{quote}
Say the Japanese for to apply and a sick visit.
\end{quote}

\subsection*{Your turn}
Say these aloud:

\begin{itemize}
\raggedright
  \item to apply, a sick visit, recovery, all right, endurance
  \item dangerous, danger, to be alive, to rescue, to suffer
\end{itemize}

\begin{tcolorbox}[breakable,colback=teal!4,colframe=teal!35!black,title={Before you move on}]
To apply? (\textbf{\ja{ぬる}}, \emph{nuru}.) A sick visit? (\textbf{\ja{おみまい}}, \emph{omimai}.) Recovery? (\textbf{\ja{かいふく}}, \emph{kaifuku}.) All right? (\textbf{\ja{だいじょうぶ}}, \emph{daijōbu}.) Endurance? (\textbf{\ja{がまん}}, \emph{gaman}.) Dangerous? (\textbf{\ja{あぶない}}, \emph{abunai}.) Danger? (\textbf{\ja{きけん}}, \emph{kiken}.) To be alive? (\textbf{\ja{いきる}}, \emph{ikiru}.) To rescue? (\textbf{\ja{たすける}}, \emph{tasukeru}.) To suffer? (\textbf{\ja{くるしむ}}, \emph{kurushimu}.)
\end{tcolorbox}
\section[(dialogue)]{Danger and Staying Alive — once more}

\label{lesson:JA-R223-review-final-223}



\begin{quote}
Say the Japanese for dangerous and danger.
\end{quote}

\subsection*{Your turn}
\begin{itemize}
\raggedright
  \item \emph{Say it:} dangerous, danger, to be alive, to rescue, to suffer
\end{itemize}

\begin{tcolorbox}[breakable,colback=teal!4,colframe=teal!35!black,title={Before you move on}]
Dangerous? (\textbf{\ja{あぶない}}, \emph{abunai}.) Danger? (\textbf{\ja{きけん}}, \emph{kiken}.) To be alive? (\textbf{\ja{いきる}}, \emph{ikiru}.) To rescue? (\textbf{\ja{たすける}}, \emph{tasukeru}.) To suffer? (\textbf{\ja{くるしむ}}, \emph{kurushimu}.)
\end{tcolorbox}
